\documentclass{rrxiv}
\rrxivid{rrxiv:2605.00008}
\rrxivversion{v1}
\rrxivprotocolversion{0.1.0}
\rrxivlicense{CC-BY-4.0}
\rrxivtopics{cs.DL,cs.IR}

\title{Many small claims, all under active replication}
\author{Example Author A}
\date{2026-05-17}

\begin{document}
\maketitle

\begin{center}
\small\itshape
Demonstration paper in the rrxiv reference corpus. The canonical machine-readable version lives at \href{https://rrxiv.com/papers/rrxiv:2605.00008}{rrxiv.com/papers/rrxiv:2605.00008}.
\end{center}

\begin{abstract}
A survey paper registering many small empirical claims about preprint discoverability. Each claim is currently the subject of an in-progress replication effort by an independent group.
\end{abstract}

\section{Introduction}
A survey paper registering many small empirical claims about preprint discoverability. Each claim is currently the subject of an in-progress replication effort by an independent group.

This document is a structured encoding of the paper in the \texttt{rrxiv} protocol's Canonical Intermediate Representation (CIR). It engages with the topics \texttt{cs.DL} and \texttt{cs.IR}. The encoding registers 5 formal claims (5 untested). Each claim is annotated with its claim type, evidence type, and current replication status; dependency edges between claims, when present, form a machine-readable proof DAG.

\section{Methodology}
We follow the \texttt{rrxiv} convention of separating \emph{claims} (the proposition under consideration) from \emph{evidence} (the argument or data supporting it). Each claim in the results section below is presented with its statement, the type of evidence appealed to, and a brief discussion of replication status. Where claims depend on prior results --- internal or external --- the dependency is recorded in the CIR as a \texttt{\textbackslash dependson} edge, so the full inferential structure is machine-traversable. Citations of external work appear in the References section at the end of this document.

\section{Results: registered claims}
\subsection*{Claim 1}
\begin{claim}[Claim 1]
\label{claim:c1}
Preprint titles longer than 12 words receive 18\% less cross-domain attention.

\emph{Replication status: untested.}
\end{claim}
This claim is an empirical observation supported by data. As of the encoding date, it has not yet been independently tested.

\subsection*{Claim 2}
\begin{claim}[Claim 2]
\label{claim:c2}
Adding a structured abstract correlates with 22\% higher click-through from search.

\emph{Replication status: untested.}
\end{claim}
This claim is an empirical observation supported by data. As of the encoding date, it has not yet been independently tested.

\subsection*{Claim 3}
\begin{claim}[Claim 3]
\label{claim:c3}
Domain-experts cite within their own subfield 4x more than cross-domain.

\emph{Replication status: untested.}
\end{claim}
This claim is an empirical observation supported by data. As of the encoding date, it has not yet been independently tested.

\subsection*{Claim 4}
\begin{claim}[Claim 4]
\label{claim:c4}
Section-level retrieval beats whole-paper retrieval on recall@5 for narrow queries.

\emph{Replication status: untested.}
\end{claim}
This claim is an empirical observation supported by data. As of the encoding date, it has not yet been independently tested.

\subsection*{Claim 5}
\begin{claim}[Claim 5]
\label{claim:c5}
The reproducibility-budget signal is stable across three independent reannotation rounds.

\emph{Replication status: untested.}
\end{claim}
This claim is an empirical observation supported by data. As of the encoding date, it has not yet been independently tested.

\section{Discussion}
The claim graph above is the primary product of this paper. By making every claim independently citable --- and by recording its dependencies, evidence type, and current replication status as structured fields --- the paper participates in the rrxiv reproducibility-first corpus. Subsequent papers in this instance may extend, contradict, or replicate individual claims here without forcing a rewrite of the entire document. See the canonical version online for the live discourse layer.

\section{References}
\begin{itemize}[leftmargin=*]
\item Discoverability metrics for preprints
\item Cross-domain attention in scholarly networks
\end{itemize}
\end{document}
